\documentclass[12pt]{article}
\usepackage{hyperref}

\usepackage{graphicx}

\usepackage{mathtools}
\DeclarePairedDelimiter\ceil{\lceil}{\rceil}
\DeclarePairedDelimiter\floor{\lfloor}{\rfloor}


\usepackage{amsmath}
\usepackage{amsfonts}
\usepackage{amssymb}

\usepackage{listings}
\lstset{language=C++,
                basicstyle=\ttfamily,
                keywordstyle=\color{blue}\ttfamily,
                stringstyle=\color{red}\ttfamily,
                commentstyle=\color{green}\ttfamily,
                morecomment=[l][\color{magenta}]{\#}
}

\usepackage{breqn}
\usepackage[]{algorithm2e}
\usepackage{physics}

\newcommand\fnurl[2]{%
  \href{#2}{#1}\footnote{\url{#2}}%
}

\usepackage{cancel}

\newcommand{\Four}{\mathcal{F}}
\renewcommand{\(}{\left(}
\renewcommand{\)}{\right)}
%% \renewcommand{\{}{\left{}
%% \renewcommand{\}}{\right}}


\title{Pruned FFTs and rocFFT}
\author{The rocFFT Team}



\def\no#1{\nobreak{#1}} % stop equation breaking in dmath

\begin{document}
\maketitle

\section{Introduction}

FFTs are often used on data where a large number of the modes are
known a priori to be zero.  For example, with low-pass filtering,
convolutions, or dealing with different resolutions.  A FFT can of
course compute the transform by reading in entire data array, but it
may be advantageous to use the sparsity pattern, thus avoiding
unnecessary calculation and data transfers. Such transforms are
referred to as pruned FFTs.

The sparsity of input data follows a block format; above a certain
index, the data is zero.  As such, compressed row or column formats
aren't typically used.  For convolutions on Hermitian-symmetric data,
$1/3$ of the inputs are zero (in each dimension); other convolutions
have $1/2$ of the inputs as zero (in each dimension).  For
applications where different fields with different resolution are
compared, this ratio may be much higher, and we have discussed a
factor of 16 in each dimension.

FFTW includes a note in their documentation on pruned FFTs at
\url{https://www.fftw.org/pruned.html}, speaking about why they are
not particularly effective in one dimension, and are not implemented
in FFTW.  There is some work by the Spiral authors
(\url{https://ieeexplore.ieee.org/document/4959642} which showed some
improvement.  The Spiral work required a modification to all kernels
(ie the sparsity pattern was used at all stages), which is a
significant development investment. cuFFT does not seem to implement a
pruned FFT, nor does Intel's oneAPI math kernel library.

\section{The blessing of dimensionality}

As usual, one-dimensional and multi-dimensional FFTs offer very
different optimisation strategies.

\subsection{One-dimensional transforms}

For pruned 1D FFTs, which the speedup is generally fairly modest, and
requires a large amount of development work.  In terms of rocFFT, the
initial strategy would be to reprocess the input data so that, for
example, only the first few non-zero entries would be read, thereby
avoiding the unnecessary global memory reads on data which is known a
priori to be zero.  However, after the first pass, the sparsity
pattern will be destroyed, so multi-kernel transforms will still
require the same number of kernels and work memory.  Moreover,
out-of-place transforms will be accomplished simply by reading the
non-zero part of the input data from global memory and the treating
the rest of the transform as an in-place FFT on the output buffer.

This will slightly improve performance because we can eliminate
global-memory reads.  There is some risk that the FFT will require a
bit more work memory because the transform is done in-place (this is
particularly a concern with real/complex transforms).  On the other
hand, while the rocFFT performance may not improve that much, this
does allow the application code to avoid an extra copy from the
compressed input buffer to the zero-padded output buffer; by removing
this extra global-memory copy, we do have some opportunities to
increase application performance.  The addition of pruned transforms
to rocFFT would require a change to the public API - this is something
that we take quite seriously.  An analysis of use cases would be
required, and we would have to come up with a strategy on how to
release this feature.

\subsection{Multi-dimensional transforms}

Multi-dimensional transforms offer a simpler and more effective
optimization strategy when dealing with sparse data.  For a
2-dimensional transform of size $N\times{}N$ with $1/2$ padding, the
FFT is computed by performing first $N$ 1-D transforms of length $N$
in the $y$ direction, and the computing $N$ 1-D transforms of length
$N$ in the $x$ direction.  Since $1/2$ of the data is zero in each
dimension, $1/2$ of the transforms $y$-direction transforms are on
zero data, and can be skipped.

For three dimensions where only a fraction $a/b$ (in each dimension)
is non-zero, the first transform can skip $(1- a/b)^2$ transforms, and
the second transforms can skip $1 - a/b$.  The last dimension of the
transform will not have any simple sparsity transform to take
advantage of, but could be combined with a one-dimensional pruned
transform if this is available in the FFT library.

By skipping sub-transforms on zero data, applications can reduce
global memory use, reduce the number of needed computations, and also
reduce the amount of MPI communication.  The multi-dimensional
optimization for pruned FFTs can be implemented at the application
level without regardless of the underlying library.  As usual,
distributed data management is an important detail - the
one-dimensional batched sub-transforms need to be contiguous, and
application developers must ensure that the MPI communications produce
the expected data layout that is consumed by the FFT library.


\end{document}

% LocalWords:  FFTs priori FFT FFTW ie cuFFT oneAPI dimensionality
% LocalWords:  optimisation rocFFT MPI
